% !TEX program = lualatex
% 自動生成: mar.report.latex — 手で編集した場合は再生成で上書きされる
\documentclass[paper=a4,fontsize=10.5pt,jafontsize=10.5pt]{jlreq}

\usepackage{luatexja}
\usepackage[haranoaji]{luatexja-preset}
\usepackage{amsmath,amssymb,amsthm,mathtools}
\usepackage{bm}
\usepackage{graphicx}
\usepackage{booktabs}
\usepackage{listings}
\usepackage{xcolor}
\usepackage[margin=25mm]{geometry}
\usepackage[hidelinks]{hyperref}
\usepackage{url}

\lstset{
  basicstyle=\ttfamily\footnotesize,
  breaklines=true,
  frame=single,
  showstringspaces=false,
  numbers=left,
  numberstyle=\tiny\color{gray},
  captionpos=b,
}

\theoremstyle{plain}
\newtheorem{theorem}{定理}[section]
\newtheorem{proposition}[theorem]{命題}
\newtheorem{lemma}[theorem]{補題}
\newtheorem{corollary}[theorem]{系}
\newtheorem{conjecture}[theorem]{予想}
\theoremstyle{definition}
\newtheorem{definition}[theorem]{定義}
\newtheorem{example}[theorem]{例}
\newtheorem{remark}[theorem]{注意}

\renewcommand{\proofname}{証明}

\title{近傍独立数の平均によるハミルトン路の十分条件: Written on the Wall II 予想 194 の証人付き網羅検証と仮定クラスの分類}
\author{自動研究パイプライン \texttt{mar}（Claude Opus 5 + 検証器）}
\date{2026年07月27日}

\begin{document}
\maketitle

\begin{abstract}
DeLaViña の Graffiti.pc が生成した未解決予想 Written on the Wall II Conjecture 194 を扱う。$\ell_{\mathrm{avg}}(G)$ を、各頂点の開近傍が誘導する部分グラフの独立数の平均とするとき、この予想は連結グラフ $G$ が $\alpha(G) \le 1 + \ell_{\mathrm{avg}}(G)$ を満たせばハミルトン路をもつと主張する。本稿ではまず手計算で、(i) $\alpha(G) \le 2$ の場合に予想が成り立つこと、(ii) 仮定を満たす位数 $n \ge 2$ の連結グラフは平均次数 $\bar{d}$ について $n < (1+\bar{d})^2$ という\textbf{密度条件}を満たすこと、その系として (iii) 連結 $r$-正則グラフでは $n \le (r+1)^2 - 1$、$n \ge 5$ の木では仮定が空虚になること、(iv) 仮定を満たす非ハミルトングラフが $n \ge 4$ のどの位数にも存在し、結論をハミルトン閉路に強化できないことを示す。そのうえで、連結グラフの完全リスト ($n \le 10$)、木の完全リスト ($n \le 20$)、GENREG の連結正則グラフ、計 13,402,242 個を網羅的に検証した。反例は存在しない。$\alpha(G)$ の計算もハミルトン路の存在判定も NP 困難だが、含意の\textbf{どちら側を閉じるかをグラフごとに選ぶ}ことで、証人 1 個あたり多項式時間で確認できる形にした: 結論が成り立つなら路 1 本を、仮定が破れるならそれを破る独立集合 1 個を渡す。仮定を満たすのは 3,337,585 個 (24.9\%) で、そのうち 3,322,958 個 (99.6\%) は $\alpha \ge 3$ であり、(i) の議論では処理できない。等号 $\alpha = 1 + \ell_{\mathrm{avg}}$ は 287,005 個で成立する。
\end{abstract}

\noindent\textbf{キーワード:} graph theory、hamiltonian path、independence number、local independence、Graffiti.pc、open problem、certificate


\section{はじめに}

$G$ を位数 $n = |V(G)|$、辺数 $m$ の連結な単純グラフとする。$\alpha(G)$ を
独立数、$N(v)$ を頂点 $v$ の\textbf{開}近傍 ($v$ 自身を含まない) とし、
\[ \ell_{\mathrm{avg}}(G) \ = \ \frac{1}{n} \sum_{v \in V(G)}
   \alpha\bigl(G[N(v)]\bigr) \]
と置く。$\alpha(G[N(v)])$ は「$v$ の隣人のうち互いに隣接しないものを最大何個
取れるか」であり、$G$ が $v$ の周りでどれだけ局所的に疎かを測る量である。
以下、その総和を
\[ S(G) \ = \ \sum_{v \in V(G)} \alpha\bigl(G[N(v)]\bigr)
   \ = \ n \, \ell_{\mathrm{avg}}(G) \]
と書く。グラフが\textbf{traceable} であるとは、ハミルトン路 (全頂点を
ちょうど 1 度ずつ通る路) をもつことをいう。

\begin{conjecture}[Graffiti.pc; WOWII Conjecture 194 \cite{wowii}]\label{conj:194}
連結グラフ $G$ が $\alpha(G) \le 1 + \ell_{\mathrm{avg}}(G)$ を満たせば、
$G$ は traceable である。
\end{conjecture}

この予想は Google DeepMind の \texttt{formal-conjectures} \cite{formalconj} に
Lean 4 で形式化されており、2026 年 7 月 27 日に取得した時点で
\texttt{@[category research open]} (未解決) と分類されている。
本稿では分母を払い、\textbf{整数の}条件
\begin{equation}\label{eq:hyp}
n \, \alpha(G) \ \le \ n + S(G)
\end{equation}
の形で扱う。以下これを\textbf{仮定}と呼ぶ。

ハミルトン路の存在判定は NP 完全であり、$\alpha(G)$ の計算も NP 困難である。
それでも予想 \ref{conj:194} は「独立数と、局所的な独立数の平均」という
2 つの量だけで traceability を保証しようとしている点で非自明である。
最も近い既知の結果は Chvátal--Erdős \cite{chvatalerdos} の系
\[ \kappa(G) \ \ge \ \alpha(G) - 1 \ \implies \ G \text{ は traceable} \]
であるが、予想 \ref{conj:194} の仮定は連結度に一切触れないので、
一方から他方は従わない。第 \ref{sec:hand} 節では両者の関係を精密にし、
第 \ref{sec:check} 節の網羅検証では
\textbf{Chvátal--Erdős では処理できないグラフが何個あるか}を数える。

\section{証人つき検証という設計}\label{sec:design}

予想 \ref{conj:194} は含意である。したがってグラフごとに、
\textbf{結論が成り立つ}ことか\textbf{仮定が破れる}ことのどちらか一方を
示せばよく、どちらを示すかはグラフごとに選んでよい。本稿の証明書はこの
選択を明示的に記録する。

\begin{lemma}[モード A: 路の証人]\label{lem:path}
頂点列 $v_1, \dots, v_n$ が $V(G)$ の並べ替えであって
$v_i v_{i+1} \in E(G)$ ($1 \le i < n$) ならば、$G$ は予想 \ref{conj:194} の
反例ではない。
\end{lemma}

証人の確認は $n-1$ 回の隣接判定と重複のチェックだけで、線形時間である。
$\alpha(G)$ も $S(G)$ も計算する必要がない。

\begin{lemma}[モード B: 独立集合の証人]\label{lem:mask}
$A \subseteq V(G)$ が独立集合で $n\,|A| > n + S(G)$ を満たすならば、$G$ は
式 \eqref{eq:hyp} を満たさず、したがって予想 \ref{conj:194} の反例ではない。
\end{lemma}

\begin{proof}
$\alpha(G) \ge |A|$ より $n\,\alpha(G) \ge n\,|A| > n + S(G)$。
\end{proof}

補題 \ref{lem:mask} の確認には $S(G)$ が要るが、これは各頂点の近傍という
\emph{小さい}部分グラフの独立数の和なので、検証器が厳密に計算し直せる
(近傍の位数は最大次数以下である)。全体の独立数 $\alpha(G)$ を解き直す必要は
ない点が重要である。

\subsection{どちらのモードを使うかを決める規則}

探索器は、まず\textbf{仮定が破れるか}を厳密に判定する。$k$ を整数として
$\alpha(G) \ge k$ は「大きさ $k$ の独立集合が存在する」ことと同値だから、
\[ k_{\min} \ = \ \left\lfloor \frac{n + S(G)}{n} \right\rfloor + 1 \]
と置けば、\emph{大きさ $k_{\min}$ の独立集合が存在すること}が
\emph{式 \eqref{eq:hyp} が破れること}と同値になる
($n\alpha > n+S \iff \alpha > \lfloor (n+S)/n \rfloor$)。
存在すればその 1 つを証人にして\textbf{モード B}、存在しなければ仮定が
成り立つので、ハミルトン路を探して\textbf{モード A} とする。
どちらの分岐でも $\alpha(G)$ を求めずに済むのが要点である。

証人は族ごとに 1 つの gzip 圧縮バイナリにまとめる。先頭に列挙順のモード
ビット列 (0 が路、1 が独立集合)、続いてモード A の路を 1 頂点
$\lceil \log_2 n \rceil$ ビットで詰めた列、最後にモード B のマスクを
$\lceil n/8 \rceil$ バイトずつ並べる。ファイルの SHA-256 を証明書 JSON に
記録し、検証器は同じ順序で元データを走査しながら証人を消費する。
独立集合の証人は\textbf{大きさ $k_{\min}$ の辞書式最小のもの}に取る。
これは「仮定が破れるか」の厳密な判定そのものであり、かつ証人が最小サイズなので
族の中で似たマスクが並び、圧縮が効く。

\subsection{検証器}

検証器 (\texttt{mar.checkgraph}) は探索器のコードを一切参照せず、標準ライブラリ
のみで書かれている。アルゴリズムも意図的に変えてある。

\begin{enumerate}
\item 近傍の独立数は\emph{極大独立集合の全列挙}の最大サイズ
      (探索器は補グラフの最大クリークを分枝限定で解く)。
\item $\alpha(G)$ も同じ全列挙による (探索器は貪欲な枝刈り付き分枝限定)。
\item ハミルトン路の\emph{確認}は隣接関係の逐次照合のみ。反例を主張された
      ときだけ Held--Karp 型のビット DP で存在判定をやり直す
      (探索器は連結性で枝刈りした深さ優先探索)。
\item 元リストの個数を照合する公表値 (OEIS A001349, A000055, A002851,
      A006820--A006822) も、探索器の表ではなく検証器が独自にもつ定数を使う。
\end{enumerate}

\subsection{どこまでを検証済みと呼ぶか}\label{sec:tier}

補題 \ref{lem:path}, \ref{lem:mask} が閉じるのは\textbf{予想の成否}だけで
ある。「仮定を満たすグラフがちょうど何個あるか」という\emph{分類}の主張は、
モード A のグラフで $\alpha(G)$ を厳密に求めない限り閉じない。そこで
本稿は族を 2 つのティアに分ける。

\begin{itemize}
\item \textbf{分類する族} (連結グラフ全体と正則グラフ): 検証器がモード A の
      グラフすべてで $\alpha(G)$ を厳密に計算し直し、式 \eqref{eq:hyp} が
      実際に成り立つこと、モード B のグラフでは実際に破れることを確かめる。
      等号の個数と例もここで閉じる。
\item \textbf{分類しない族} (木): 系 \ref{cor:tree} により $n \ge 5$ の木は
      仮定を満たさないので、全グラフがモード B になる。この族について本稿が
      主張するのは「仮定が空虚である」ことだけであり、これは
      補題 \ref{lem:mask} だけで閉じる。
\end{itemize}

\section{手計算でどこまで詰められるか}\label{sec:hand}

網羅検証に入る前に、予想 \ref{conj:194} がどこまで\emph{証明できる}か、
また仮定がどれだけ強いかを見ておく。以下 $G$ は位数 $n \ge 2$ の連結グラフ、
$\bar{d} = 2m/n$ を平均次数とする。

\begin{lemma}\label{lem:range}
$1 \le \ell_{\mathrm{avg}}(G) \le \alpha(G)$ である。したがって仮定
\eqref{eq:hyp} は $\alpha(G) - 1 \le \ell_{\mathrm{avg}}(G) \le \alpha(G)$
と同値であり、「局所的な独立数の平均が全体の独立数にほぼ等しい」ことを
要求している。
\end{lemma}

\begin{proof}
連結かつ $n \ge 2$ なので各 $v$ で $N(v) \ne \emptyset$、よって
$\alpha(G[N(v)]) \ge 1$。また $G[N(v)]$ の独立集合は $G$ の独立集合でもあるから
$\alpha(G[N(v)]) \le \alpha(G)$。平均を取れば主張を得る。
\end{proof}

\begin{theorem}\label{thm:alpha2}
$\alpha(G) \le 2$ ならば $G$ は仮定 \eqref{eq:hyp} を満たし、かつ traceable
である。したがって予想 \ref{conj:194} は $\alpha(G) \le 2$ の範囲で成立し、
反例は $\alpha(G) \ge 3$ を満たす。
\end{theorem}

\begin{proof}
補題 \ref{lem:range} より $\ell_{\mathrm{avg}} \ge 1$ なので
$\alpha \le 2 \le 1 + \ell_{\mathrm{avg}}$、すなわち仮定は自動的に成り立つ。
一方 $G$ は連結なので $\kappa(G) \ge 1 \ge \alpha(G) - 1$ であり、
Chvátal--Erdős \cite{chvatalerdos} の系より traceable である。
\end{proof}

定理 \ref{thm:alpha2} は予想 \ref{conj:194} の一部を証明すると同時に、
\textbf{仮定を満たすグラフのうち $\alpha \le 2$ のものは既知の結果で尽きる}
ことを意味する。網羅検証で数えるべきはその外側、$\alpha \ge 3$ の部分である。

\begin{theorem}[密度条件]\label{thm:density}
$n \ge 2$ の連結グラフ $G$ が仮定 \eqref{eq:hyp} を満たすならば
\[ n \ < \ (1 + \bar{d})^2, \qquad \text{すなわち} \qquad
   m \ > \ \frac{n(\sqrt{n} - 1)}{2} \]
である。とくに仮定を満たすグラフ族は平均次数が $\sqrt{n} - 1$ より速く
増える\emph{密な}族に限られる。
\end{theorem}

\begin{proof}
$\alpha(G[N(v)]) \le |N(v)| = \deg(v)$ より
\begin{equation}\label{eq:S2m}
S(G) \ \le \ \sum_{v} \deg(v) \ = \ 2m \ = \ n \bar{d} .
\end{equation}
一方 Caro--Wei の下界 \cite{carowei} より
$\alpha(G) \ge \sum_{v} \frac{1}{\deg(v)+1} \ge \frac{n}{\bar{d}+1}$
(後半は Cauchy--Schwarz、すなわち調和平均と算術平均の不等式)。これを仮定
\eqref{eq:hyp} に入れると
\[ \frac{n^2}{\bar{d}+1} \ \le \ n\,\alpha(G) \ \le \ n + S(G)
   \ \le \ n(1 + \bar{d}) \]
となり $n \le (1+\bar{d})^2$ を得る。等号が成り立つには
\eqref{eq:S2m} と Caro--Wei の両方で等号が要る。前者の等号は各近傍が
独立集合であること (すなわち $G$ が三角形をもたないこと)、後者の等号は
$G$ がクリークの非交和であることと同値で、さらに調和平均と算術平均が
一致するには $G$ が正則でなければならない。三角形をもたないクリークの非交和は
$K_1$ と $K_2$ だけからなるので、連結性と合わせると $G = K_1$ または
$G = K_2$ である。$G = K_2$ では $(1+\bar{d})^2 = 4 \neq 2 = n$ なので
等号は成り立たず、等号が起きるのは $n = 1$ に限る。よって $n \ge 2$ では
狭義の不等式が成り立つ。
\end{proof}

\begin{corollary}\label{cor:reg}
連結 $r$-正則グラフが仮定 \eqref{eq:hyp} を満たすならば
$n \le (r+1)^2 - 1$ である。とくに位数 $16$ 以上の連結 3-正則グラフ、
位数 $25$ 以上の連結 4-正則グラフでは仮定が空虚である。
\end{corollary}

\begin{proof}
定理 \ref{thm:density} で $\bar{d} = r$ とし、$n$ が整数であることを使う。
\end{proof}

\begin{corollary}\label{cor:tree}
$n \ge 5$ の木は仮定 \eqref{eq:hyp} を満たさない。すなわち予想
\ref{conj:194} は木に対して空虚に真である。
\end{corollary}

\begin{proof}
木は三角形をもたないので各近傍は独立集合であり、\eqref{eq:S2m} は等号、
$S = 2m = 2(n-1)$。よって仮定は $n\alpha \le 3n - 2$、すなわち
$\alpha \le 3 - 2/n < 3$、整数性から $\alpha \le 2$ である。一方、木は
二部グラフなので大きい側の部集合を取って $\alpha \ge \lceil n/2 \rceil$ で
あり、$n \ge 5$ では $\alpha \ge 3$ となって矛盾する。
\end{proof}

\begin{proposition}[二部グラフへの縮約]\label{prop:bip}
部集合の大きさが $x \ge y$ の連結二部グラフ $G$ が仮定 \eqref{eq:hyp} を
満たすならば $x - y \le 1$ である。さらに
\begin{itemize}
\item $x = y = k$ のとき $\alpha(G) = k$ (同値に $G$ は完全マッチングをもつ)
      であり、$K_{k,k}$ から取り除かれている辺の本数は $k$ 以下である。
\item $x = k+1$, $y = k$ のとき $K_{k+1,k}$ から取り除かれている辺の本数は
      $k/2$ 以下である。
\end{itemize}
\end{proposition}

\begin{proof}
二部グラフは三角形をもたないので \eqref{eq:S2m} は等号 $S = 2m$、
また $\alpha \ge x$、$m \le xy$ である。仮定 \eqref{eq:hyp} より
$(x+y)x \le (x+y) + 2xy$、整理して $x(x-y) \le x+y$。
$x - y \ge 2$ なら左辺 $\ge 2x \ge x+y$ で、等号には $x=y$ が要るので矛盾。
よって $x-y \le 1$。

$x=y=k$ とする。$\alpha \ge k+1$ を仮定すると $2k(k+1) \le 2k + 2m$ から
$m \ge k^2$、すなわち $G = K_{k,k}$ となるが $\alpha(K_{k,k}) = k$ で矛盾。
よって $\alpha = k$ で、König の定理より最大マッチングの大きさは
$2k - \alpha = k$、すなわち完全マッチングが存在する。また $\alpha = k$ を
仮定に入れて $2k \cdot k \le 2k + 2m$、$m \ge k^2 - k$ を得る。
$x = k+1$, $y = k$ のときは $\alpha \ge k+1$ と $n = 2k+1$ から
$(2k+1)(k+1) \le (2k+1) + 2m$、$m \ge k^2 + k/2$ である。
\end{proof}

命題 \ref{prop:bip} は、二部グラフに限れば予想 \ref{conj:194} が
「$K_{k,k}$ から高々 $k$ 本の辺を除いた連結グラフは traceable か」という
具体的な問題に縮約されることを示している。本稿ではこの縮約までを与え、
一般の $k$ に対する証明は与えない (第 \ref{sec:limit} 節)。

\begin{proposition}[仮定クラスは全位数で非空、かつ結論は最良]\label{prop:sharp}
$k \ge 3$ とし、$G_k$ を完全グラフ $K_k$ の 1 頂点に次数 1 の頂点を 1 つ
接続したグラフとする ($n = k+1$)。このとき $\alpha(G_k) = 2$,
$S(G_k) = n+1$ であって $G_k$ は仮定 \eqref{eq:hyp} を満たし、traceable だが
ハミルトン閉路をもたない。
\end{proposition}

\begin{proof}
$K_k$ の頂点を $w_1, \dots, w_k$、懸垂頂点を $u$ とし $u \sim w_1$ とする。
$u$ と $w_i$ ($i \ge 2$) は隣接しないので $\alpha \ge 2$、$K_k$ の 2 頂点は
つねに隣接し $u$ を含む独立集合は $\{u, w_i\}$ の形に限るので $\alpha = 2$。
近傍の独立数は、$w_i$ ($i \ge 2$) ではクリークなので $1$、$w_1$ では
$\{u, w_j\}$ が取れて $2$、$u$ では $1$。よって
$S = (k-1) + 2 + 1 = k + 2 = n + 1$ で、
$n\alpha = 2n \le 2n + 1 = n + S$ が成り立つ。
$u, w_1, w_2, \dots, w_k$ はハミルトン路である。$u$ の次数が $1$ なので
ハミルトン閉路は存在しない。
\end{proof}

命題 \ref{prop:sharp} から次の 2 点が従う。第一に、仮定 \eqref{eq:hyp} は
どの位数 $n \ge 4$ でも空虚ではない。第二に、予想 \ref{conj:194} の結論を
「ハミルトン閉路をもつ」に強化することはできない。もっとも $\alpha(G_k) = 2$
なので、この族は定理 \ref{thm:alpha2} の射程内にある。$\alpha \ge 3$ かつ
仮定を満たすグラフが実際に多数存在することは、次節の網羅検証が示す。

\section{網羅検証}\label{sec:check}

\begin{theorem}\label{thm:main}
位数 $n \le 10$ の連結グラフ全体 (11,989,763 個)、
位数 $11 \le n \le 20$ の木全体 (1,345,823 個。位数 $10$
以下の木は前者に含まれる)、および
\[ (n, r) \ \in \ \{(12, 3), (14, 3), (16, 3), (18, 3), (11, 4), (12, 4), (13, 4), (12, 5), (11, 6)\} \]
なる連結 $r$-正則グラフ全体 (66,656 個) のいずれにおいても、予想
\ref{conj:194} は成立する。とくに反例 $G$ が存在するならば
$n \ge 11$ であり、$G$ が正則ならば上記の族に属さない。
\end{theorem}

\begin{proof}
表 \ref{tab:fams} の各族について、グラフごとに補題 \ref{lem:path} の路か
補題 \ref{lem:mask} の独立集合のいずれかを記録した。証明書
\texttt{p0006\_wowii194\_hamiltonian.json} と証人ファイルに対して検証器がすべての証人を確認し、
全 13,402,242 個で成立した (モード A が 3,337,585 個、
モード B が 10,064,657 個)。元リストの完全性は、走査行数が OEIS A001349
(連結グラフ)、A000055 (木)、A002851 および A006820--A006822 (連結正則グラフ)
の公表値と一致することで担保される。
\end{proof}

\begin{table}[htbp]
\centering
\caption{走査した族と分類の内訳。「仮定」は式 \eqref{eq:hyp} を満たす個数、
「$\alpha \ge 3$」はそのうち定理 \ref{thm:alpha2} の射程外にあるもの、
「等号」は $\alpha = 1 + \ell_{\mathrm{avg}}$ となるもの、
「証人 B」は仮定が破れることを独立集合で示した個数である。等号欄の「---」は
分類しない族 (第 \ref{sec:tier} 節) で、0 個という意味ではない。}
\label{tab:fams}
\begin{tabular}{lrrrrrr}
\toprule
族 & $n$ & 個数 & 仮定 & $\alpha \ge 3$ & 等号 & 証人 B \\
\midrule
連結グラフ & 2 & 1 & 1 & 0 & 0 & 0 \\
連結グラフ & 3 & 2 & 2 & 0 & 0 & 0 \\
連結グラフ & 4 & 6 & 5 & 0 & 0 & 1 \\
連結グラフ & 5 & 21 & 14 & 2 & 1 & 7 \\
連結グラフ & 6 & 112 & 62 & 27 & 13 & 50 \\
連結グラフ & 7 & 853 & 417 & 313 & 107 & 436 \\
連結グラフ & 8 & 11,117 & 5,120 & 4,714 & 1,025 & 5,997 \\
連結グラフ & 9 & 261,080 & 98,994 & 97,101 & 10,576 & 162,086 \\
連結グラフ & 10 & 11,716,571 & 3,226,062 & 3,213,895 & 273,633 & 8,490,509 \\
木 & 11 & 235 & 0 & 0 & --- & 235 \\
木 & 12 & 551 & 0 & 0 & --- & 551 \\
木 & 13 & 1,301 & 0 & 0 & --- & 1,301 \\
木 & 14 & 3,159 & 0 & 0 & --- & 3,159 \\
木 & 15 & 7,741 & 0 & 0 & --- & 7,741 \\
木 & 16 & 19,320 & 0 & 0 & --- & 19,320 \\
木 & 17 & 48,629 & 0 & 0 & --- & 48,629 \\
木 & 18 & 123,867 & 0 & 0 & --- & 123,867 \\
木 & 19 & 317,955 & 0 & 0 & --- & 317,955 \\
木 & 20 & 823,065 & 0 & 0 & --- & 823,065 \\
3-正則 & 12 & 85 & 0 & 0 & 0 & 85 \\
3-正則 & 14 & 509 & 0 & 0 & 0 & 509 \\
3-正則 & 16 & 4,060 & 0 & 0 & 0 & 4,060 \\
3-正則 & 18 & 41,301 & 0 & 0 & 0 & 41,301 \\
4-正則 & 11 & 265 & 96 & 96 & 25 & 169 \\
4-正則 & 12 & 1,544 & 353 & 353 & 147 & 1,191 \\
4-正則 & 13 & 10,778 & 289 & 289 & 136 & 10,489 \\
5-正則 & 12 & 7,848 & 5,905 & 5,905 & 1,342 & 1,943 \\
6-正則 & 11 & 266 & 265 & 263 & 0 & 1 \\
\midrule
\multicolumn{2}{l}{合計} & 13,402,242 & 3,337,585 & 3,322,958 &
287,005 & 10,064,657 \\
\bottomrule
\end{tabular}
\end{table}

\begin{corollary}\label{cor:beyond}
走査した 13,402,242 個のうち仮定 \eqref{eq:hyp} を満たすのは 3,337,585 個
(24.9\%) であり、そのうち 3,322,958 個 (99.6\%) は
$\alpha(G) \ge 3$ を満たす。すなわちこれらは定理 \ref{thm:alpha2} の
射程外にあり、予想 \ref{conj:194} の内容が実際に問われるのはこの部分である。
\end{corollary}

\begin{proof}
分類する族では検証器が $\alpha(G)$ を厳密に計算し直し、モード A のグラフで
式 \eqref{eq:hyp} が成り立つこと、モード B のグラフで破れることを確かめた
(第 \ref{sec:tier} 節)。$\alpha \ge 3$ の判定は、分類しない族でも
極大独立集合の列挙から独立に行っている。
\end{proof}

ただし $\alpha(G) \ge 3$ であっても、連結度が高ければ Chvátal--Erdős の系
($\kappa \ge \alpha - 1$ ならば traceable) が結論を与えることがある。本稿は
$\kappa(G)$ を計算していないので、上の 3,322,958 個は「既知の十分条件では
処理できないグラフ」の個数そのものではなく、その\textbf{上界}である。

系 \ref{cor:tree} が予言するとおり、走査した木の族 ($11 \le n \le 20$、1,345,823 個) では仮定を満たすグラフが 1 個も現れなかった (表 \ref{tab:fams})。この族の寄与は、予想が真であることの証拠ではなく、\textbf{仮定が空虚である}ことの独立な確認である。

系 \ref{cor:reg} により、$(16, 3)$・$(18, 3)$ の族では仮定が空虚であり、実際に仮定成立グラフは 0 個である。

\begin{corollary}\label{cor:eq}
予想 \ref{conj:194} の仮定は本稿の範囲で最良である。走査した族のうち分類を
行ったもの (12,056,419 個) の中で、等号
$\alpha(G) = 1 + \ell_{\mathrm{avg}}(G)$ を達成するグラフは 287,005 個
存在し、いずれも traceable である。ただし位数 $10$ の連結グラフ (273,633 個) については、等号グラフの個数と例は検証済みだが、上限 200,000 を超えるため全リストは証明書に含めていない。
\end{corollary}

等号を達成するグラフを graph6 表記で挙げる。位数が最小のものは位数 $5$ の \texttt{DEw}（この位数では 1 個に限る）である。走査した最大の位数 $10$ では 273,633 個あり、たとえば \texttt{I?ABvrw\textasciicircum{}?}、\texttt{I?ABvrw\}?}、\texttt{I?ABvr\{\textasciitilde{}?} である。等号は正則グラフでも起こる。最小の例は $4$-正則・位数 $11$ の \texttt{J\textasciitilde{}ooOSE@WF\_} である。

\section{限界}\label{sec:limit}

本稿が示したのは第 \ref{sec:hand} 節の部分的な証明と、有限範囲での網羅検証
であって、予想 \ref{conj:194} の完全な証明ではない。連結グラフの完全リストは
McKay \cite{mckay} が $n \le 10$ まで公開しており、本稿はその全体と、
木については $n \le 20$ を走査した。正則グラフは
GENREG \cite{genreg} が公開している範囲から、連結グラフの族に含まれない
$n \ge 11$ のものを選んだ。

定理 \ref{thm:density} により、反例は\textbf{密}でなければならない
($m > n(\sqrt{n}-1)/2$)。密なグラフは traceable でありやすいので、これは
予想が真であることを支持する状況証拠だが、証明ではない。実際、密度条件は
$n$ が大きくなるほど相対的には弱くなる ($\bar{d}/n \to 0$ でも
$\bar{d} > \sqrt{n}-1$ は満たせる) ので、この方向だけで一般の証明は
得られない。

未解決として残るのは、定理 \ref{thm:alpha2} の外側、すなわち
$\alpha(G) \ge 3$ かつ仮定 \eqref{eq:hyp} を満たすグラフである。
本稿の走査ではそれが 3,322,958 個現れ、すべて traceable であった。
二部グラフに限れば命題 \ref{prop:bip} により問題は「$K_{k,k}$ から高々
$k$ 本の辺を除いた連結グラフの traceability」に縮約されるが、一般の $k$ に
対する証明は本稿では与えていない ($n = 2k$ なので、走査した範囲は
$k \le 5$ に相当する)。

もう 1 つ残るのは、$\alpha(G) \ge 3$ かつ仮定を満たすグラフのうち、
Chvátal--Erdős の系 ($\kappa \ge \alpha - 1$) でも処理できないものを
切り分けることである。連結度 $\kappa$ の計算は $\alpha$ が小さい範囲では
安価 (大きさ $\alpha - 2$ の頂点集合をすべて除いて連結性を見ればよい) なので、
同じ証人の枠組みに載せられる。本稿ではこれを行っていないため、系
\ref{cor:beyond} の 3,322,958 個は上界にとどまる。

Graffiti.pc 自身が開発時に小さいグラフのデータベース上で予想を試している
可能性が高い。したがって「位数の小さい連結グラフに反例がない」ことの新規性は
限定的であり、本稿の寄与は (i) 含意の両側を 1 個の証人で閉じる形に落とし、
第三者が独立に再実行できる証明書にしたこと、(ii) 仮定クラスと等号グラフを
位数 $10$ まで完全に分類し、既知の結果 (定理 \ref{thm:alpha2}) で
処理できない部分を 3,322,958 個と定量したこと、(iii) 第 \ref{sec:hand} 節の
密度条件により仮定クラスの形を決定したこと、の 3 点にある。


\section*{再現性}
本稿の主張はすべて、探索器とは独立に実装された検証器によって再検査されている。次のコマンドで再現できる。

\begin{center}\texttt{python -m mar verify p0006\_wowii194\_hamiltonian}\end{center}

\begin{center}\begin{tabular}{ll}\toprule
証明書ダイジェスト & \texttt{26e70873f0d05e73} \\
生成日時 (UTC) & 2026-07-26T19:28:34+00:00 \\
Python & 3.10.11 \\
実行環境 & Windows 11 (build 26200) \\
リポジトリ版 & \texttt{703ee1a} \\
乱数種 & 0 \\
探索時間 & 1070.6 秒 \\
\bottomrule\end{tabular}\end{center}


\begin{thebibliography}{99}
\bibitem{wowii} E. DeLaViña, Written on the Wall II: Conjectures of Graffiti.pc, University of Houston--Downtown. \url{http://cms.dt.uh.edu/faculty/delavinae/research/wowII/}
\bibitem{formalconj} Google DeepMind, formal-conjectures: FormalConjectures/WrittenOnTheWallII/GraphConjecture194.lean (2026-07-27 取得). \url{https://github.com/google-deepmind/formal-conjectures}
\bibitem{chvatalerdos} V. Chvátal, P. Erdős, A note on Hamiltonian circuits, Discrete Math. 2 (1972) 111--113.
\bibitem{carowei} Y. Caro, New results on the independence number, Tech. Report, Tel-Aviv Univ. (1979); V. K. Wei, A lower bound on the stability number of a simple graph, Bell Lab. Tech. Memo. 81-11217-9 (1981).
\bibitem{mckay} B. D. McKay, A. Piperno, Practical graph isomorphism II, J. Symbolic Comput. 60 (2014) 94--112. データ: Combinatorial Data. \url{https://users.cecs.anu.edu.au/~bdm/data/graphs.html}
\bibitem{genreg} M. Meringer, Fast generation of regular graphs and construction of cages, J. Graph Theory 30 (1999) 137--146. \url{https://www.mathe2.uni-bayreuth.de/markus/reggraphs.html}
\end{thebibliography}

\end{document}
